% =====================================================================
% Engineered Intermediate-Mass Black Hole Systems: Infrastructure
% Constraints, Observable Residue, and a Multi-Messenger Adjudication
% Framework, with Omega Centauri as Worked Example
% Tim Swanson — The Omega Centauri Society / Post Oak Labs
% Paper E. Companions: mth-paper (A), inward-review (B),
% campaign-paper (C), economics-paper (D).
% Build: pdflatex engineered-imbh-paper && bibtex engineered-imbh-paper
%        && pdflatex x2
% Bibliography: references.bib (shared) + engineered-extra.bib (new).
% =====================================================================
\documentclass[11pt]{article}

\usepackage[margin=1.1in]{geometry}
\usepackage{newtxtext,newtxmath}
\usepackage[round,authoryear]{natbib}
\usepackage{graphicx}
\usepackage{booktabs}
\usepackage{amsmath}
\usepackage[font=small,labelfont=bf]{caption}
\usepackage{microtype}
\usepackage[colorlinks=true,linkcolor=blue!50!black,citecolor=blue!50!black,urlcolor=blue!50!black]{hyperref}
\usepackage{xcolor}
\usepackage[obeyFinal,textsize=small]{todonotes}

\newcommand{\msun}{\ensuremath{M_\odot}}
\newcommand{\astar}{\ensuremath{a_\star}}
\newcommand{\ocen}{\ensuremath{\omega}~Cen}
\newcommand{\rg}{\ensuremath{r_{\mathrm{g}}}}
\newcommand{\lnK}{\ensuremath{\ln K}}
\newcommand{\Heng}{\ensuremath{H_{\mathrm{eng}}}}
\newcommand{\Hgas}{\ensuremath{H_{\mathrm{q}}}}
\newcommand{\Hsub}{\ensuremath{H_{\mathrm{sub}}}}
\newcommand{\specbox}[1]{\par\smallskip\noindent\fbox{\parbox{0.97\linewidth}{\small\textbf{Epistemic status:} #1}}\smallskip\par}

\title{\textbf{Engineered Intermediate-Mass Black Hole Systems:\\
Infrastructure Constraints, Observable Residue,\\
and a Multi-Messenger Adjudication Framework}}
\author{Tim Swanson\\[2pt]
\small The Omega Centauri Society / Post Oak Labs\\
\small \texttt{tim@postoaklabs.com}}
\date{July 2026 \\[4pt] \small Draft v0.3 (last revised 2026-07-18) --- fifth paper of the set; companion to \emph{The Macro Transcension Hypothesis} (Paper A),\\ \small the inward-migration review (Paper B), the \ocen{} campaign (Paper C), and the migration economics (Paper D);\\ \small prepared for omegacentauri.me}

\begin{document}
\maketitle

\begin{abstract}
\noindent
The inward-migration resolutions of the Fermi paradox identify rapidly spinning intermediate-mass black holes (IMBHs) in dense, old stellar clusters as thermodynamically privileged destinations for computation-optimizing civilizations. Previous work in this series argued the thermodynamic case (Paper A), surveyed the hypothesis family (Paper B), designed a multi-messenger observational campaign for the nearest candidate system, Omega Centauri (Paper C), and priced the migration decision (Paper D). Two questions remain open, and this paper addresses both. First, \emph{feasibility}: no published analysis asks whether large-scale computational infrastructure can persist in the actual environment such a destination provides, a Kerr IMBH embedded in a cluster core with stellar density $\sim 3\times10^{3}\,\msun\,\mathrm{pc}^{-3}$ and velocity dispersion $\sim 20\,\mathrm{km\,s^{-1}}$. Extending recent passive-stability results for stellar engines and Dyson bubbles \citep{McInnes2026} to the combined Kerr-plus-cluster potential, and combining analytic tidal, thermal, and material limits with a Monte Carlo of gravitationally focused stellar flybys, we derive an allowed envelope in the (orbital radius, architecture) plane for a fiducial $2\times10^{4}\,\msun$ hole in the \ocen{} core: precession-tolerant swarms survive passively from $\sim 10^{2}$ gravitational radii out to the cluster stripping radius at $\sim 4\times10^{3}$~AU, and, contrary to our initial expectation, stellar flybys never bind: the Monte Carlo yields a scale-free, conservative diffusion floor of $\sim 3\times10^{9}$~yr throughout the envelope, further lengthened at depth by adiabatic protection. Tidal disruption events set the hazard-recurrence horizon ($\gtrsim 10^{7}$~yr), black-hole Brownian wander is benign for the installation while biasing the observer's astrometry, and the measured intracluster medium ($n_{e} \simeq 0.23\,\mathrm{cm^{-3}}$) funds, via Bondi accretion at magnetically-arrested efficiencies, an estimated power budget of $\sim 10^{6}\,L_{\odot}$: supply never binds; thermal concealment does. We further show, extending the collisional-cascade analysis of \citet{Lacki2025Cascade}, that an \emph{abandoned} deep swarm grinds to debris on an estimated $10^{2}$--$10^{3}$~yr timescale and is then drained by the hole, leaving spin, which Blandford--Znajek use erodes only on $\gg$~Hubble timescales at ambient fueling, as the only durable fossil of engineered history. Second, \emph{residue and adjudication}: the feasibility envelope implies two forward-modeled observables, a Landauer-limited waste-heat floor and a magnetically-arrested-disk (MAD) regulation signature, namely suppression of the flux-eruption variability that general-relativistic magnetohydrodynamic simulations identify as characteristic of natural MAD accretion \citep{Universe2026MAD}. We then construct a hierarchical Bayesian framework, with per-messenger Bayes factors against an explicit menu of astrophysical null hypotheses (gas-starved quiescent IMBH; no IMBH, with a stellar-remnant subcluster) combined through coincidence likelihoods of the kind developed for gravitational-wave counterpart searches \citep{Ashton2018,Breschi2024}, and with pre-registered decision thresholds. Applied end-to-end to the current \ocen{} data (fast-star kinematics, pulsar-timing bounds, deep radio and infrared silence), the framework yields posterior odds that favor the astrophysical nulls, as expected, while quantifying for the first time how much each planned observation from Paper C can move the odds. The framework is stated instrument-agnostically and is usable by any technosignature program. All results are conditional on the optimization premise shared by the hypothesis family; the contribution is the feasibility envelope, the forward-modeled residue, and the adjudication machinery.
\end{abstract}

\bigskip
\noindent\textbf{Keywords:} Fermi paradox; SETI; technosignatures; intermediate-mass black holes; Omega Centauri; megastructures; magnetically arrested disks; Bayesian model selection; multi-messenger astronomy

\newpage
\tableofcontents
\newpage

% =====================================================================
\section{Introduction}
\label{sec:intro}

\subsection{The missing quadrant}

Four papers in this series have examined the hypothesis that computation-optimizing civilizations migrate toward rapidly spinning massive black holes in dense old stellar systems. Paper A \citep{Swanson2026MTH} argued \emph{why}: the thermodynamic advantages of Kerr accretion efficiency, horizon entropy disposal, and Bekenstein-bounded information storage. Paper B \citep{Swanson2026Review} surveyed \emph{who}: the six-member family of inward-migration proposals and their falsifiability grades. Paper C \citep{Swanson2026Campaign} planned \emph{how to look}: eight instrument-matched programs targeting Omega Centauri (NGC~5139), the nearest strong IMBH candidate. Paper D \citep{Swanson2026Economics} priced \emph{whether to go}: closed-form crossover conditions between migration, densification, and seeding strategies.

Two questions remain unaddressed, in this series and, so far as we can determine, in the literature. First: could the infrastructure such a civilization requires actually \emph{persist} at the destination? The destination is a demanding one. A cluster-core IMBH sits in a stellar environment three to four orders of magnitude denser than the solar neighborhood, threaded by gravitationally focused stellar traffic, and (if fueled) surrounded by a relativistic accretion flow. Black-hole computing proposals \citep{Inoue2011,Opatrny2017,Dvali2023} treat the hole as an idealized resource and the environment as empty. Megastructure stability analyses \citep{Wright2020,Raval2024,McInnes2026} treat single stars in isolation. The intersection, engineered structures bound to a Kerr hole inside a live cluster core, is unexamined.

Second: if the campaign of Paper C ever records an anomaly, by what procedure would the community decide what it had seen? Radio SETI has mature per-signal verification practice \citep{Sheikh2021,Luan2026} and Bayesian population inference \citep{Grimaldi2018}, and multi-messenger astrophysics has quantitative coincidence formalisms \citep{Ashton2018,Breschi2024}, but no published framework adjudicates a candidate \emph{technosignature} across multiple messengers against an explicit menu of astrophysical alternatives. Paper C's rule that ``every anomaly is adjudicated by at least two independent messengers'' was stated as policy; here we give it mathematical content.

The two questions are coupled, and that coupling is the reason they share a paper. A feasibility envelope is a prior: it tells the adjudicator where in parameter space an engineered system could sit, and therefore which anomalies deserve elevated scrutiny and which are excluded on engineering grounds regardless of how anomalous they appear. Forward-modeled residue channels are likelihoods: they specify what an engineered system \emph{would} look like, channel by channel, so that Bayes factors can be computed rather than gestured at. The paper thus runs in one direction: constraints (\S\ref{sec:envelope}) produce observables (\S\ref{sec:residue}) which feed an adjudication engine (\S\ref{sec:adjudication}) that we exercise on real data (\S\ref{sec:worked}).

\subsection{Conditionality and register}

\specbox{Everything downstream of \S\ref{sec:envelope} is conditional on the optimization premise examined critically in Papers A and B: that some long-lived technological lineages maximize discounted computation and act on the thermodynamic gradients described there. This paper does not defend that premise; it asks what follows from it observationally, and how a null or a detection would be established. The adjudication framework of \S\ref{sec:adjudication} is independent of the premise and applicable to any technosignature search.}

We work throughout with the fiducial system of Papers A and C: a black hole of mass $M = 2\times10^{4}\,\msun$ (bracketed by the $8.2\times10^{3}\,\msun$ kinematic lower bound of \citealt{Haberle2024Nature} and the $\sim 5\times10^{4}\,\msun$ $N$-body preferred value of \citealt{GonzalezPrieto2025}), spin left free, embedded in the \ocen{} core: central density $\rho_{0} \simeq 3\times10^{3}\,\msun\,\mathrm{pc}^{-3}$, core radius $r_{c} \simeq 3.6$~pc, line-of-sight velocity dispersion $\sigma \simeq 20$--$23\,\mathrm{km\,s^{-1}}$ near the center \citep{BaumgardtVasiliev2021,Nitschai2023}. Where a single value is needed we adopt $\sigma = 21\,\mathrm{km\,s^{-1}}$, the midpoint of that central range, and use it consistently in the influence radius of \S\ref{sec:outer}, the flyby velocities of \S\ref{sec:mc}, and Appendix~\ref{app:mc}. The series' public calculators default instead to the cluster-averaged $18.2\,\mathrm{km\,s^{-1}}$ of the shared measurement compilation, which is a global rather than central-region figure; the difference propagates as a factor $(21/18.2)^{2} \simeq 1.3$ in $r_{\mathrm{infl}}$ and is within the parameter variations of Appendix~\ref{app:mc}.

The gravitational radius is
\begin{equation}
\rg \equiv \frac{GM}{c^{2}} \simeq 3.0\times10^{9}\,\mathrm{cm}\,
\left(\frac{M}{2\times10^{4}\,\msun}\right)
\simeq 2.0\times10^{-4}\,\mathrm{AU}\left(\frac{M}{2\times10^{4}\,\msun}\right),
\end{equation}
so the dynamic range between horizon scale and cluster scale is nine orders of magnitude in radius. The feasibility question is where in that range infrastructure can live.

% =====================================================================
\section{The feasibility envelope}
\label{sec:envelope}

\subsection{Constraint inventory}
\label{sec:inventory}

We consider a one-parameter family of architectures indexed by compactness: at one end, a dense swarm of independent elements on near-circular orbits at $r \sim 10^{2}$--$10^{4}\,\rg$ (the configuration favored by the thermodynamic argument of Paper A, since proximity to the horizon minimizes entropy-transport losses); at the other, an extended bubble or shell of elements at $r \sim 10$--$10^{3}$~AU supported partly by radiation pressure or tether stress, the black-hole analogue of the Dyson bubbles whose collective stability was recently established by \citet{McInnes2026}. Intermediate cases (rings, nested tori) inherit constraints from both ends \citep{Raval2024}. We deliberately do not commit to an architecture; the envelope is the set of $(r, \text{architecture})$ pairs surviving all constraints simultaneously.

Six constraints bound the envelope:
(i) relativistic orbit stability (inner);
(ii) tidal stress on extended elements (inner);
(iii) accretion-flow radiation and magnetic environment, if the hole is fueled (inner);
(iv) cluster tidal truncation of bound orbits (outer);
(v) stellar flyby perturbation and collision hazard (outer);
(vi) thermal rejection capacity, which couples to the residue analysis of \S\ref{sec:residue}.

\subsection{Inner boundary}
\label{sec:inner}

\textbf{Orbit stability.} Circular equatorial orbits around a Kerr hole are stable outside the innermost stable circular orbit, $r_{\mathrm{ISCO}} = 6\,\rg$ (Schwarzschild) shrinking to $\rg$ (prograde, extremal) \citep{BardeenPressTeukolsky1972}. For $M = 2\times10^{4}\,\msun$ this is $\sim 10^{-3}$~AU: dynamically, matter can orbit extraordinarily deep. Orbital periods there are $P \simeq 2\pi (r^{3}/GM)^{1/2} \sim$ minutes, and communication latency across the swarm is milliseconds, which is part of the computational attraction (Paper A, \S4).

\textbf{Tidal stress.} A rigid element of size $\ell$ at radius $r$ experiences differential acceleration $\sim 2GM\ell/r^{3}$. Requiring internal stress below a material strength $S$ for an element of density $\rho_{\mathrm{el}}$ gives
\begin{equation}
\ell \lesssim \left(\frac{S\,r^{3}}{2GM\rho_{\mathrm{el}}}\right)^{1/2}
\simeq 2\times10^{2}\,\mathrm{km}
\left(\frac{S}{10\,\mathrm{GPa}}\right)^{1/2}
\left(\frac{\rho_{\mathrm{el}}}{10^{3}\,\mathrm{kg\,m^{-3}}}\right)^{-1/2}
\left(\frac{r}{10^{2}\,\rg}\right)^{3/2}
\left(\frac{M}{2\times10^{4}\,\msun}\right)^{-1},
\end{equation}
where $10$~GPa is the tensile strength of present-day carbon composites. Tides therefore never forbid the swarm architecture (elements of km scale are unconstrained down to $\sim 10\,\rg$); they forbid \emph{monolithic} structures below $\sim 10^{3}\,\rg$ and thereby select swarms at small radii, a conclusion that parallels the single-star case \citep{Wright2020}.

\textbf{Radiation and magnetic environment.} If the hole is fueled at the level required for the Blandford--Znajek power budget of Paper A, the inner $\sim 10^{2}\,\rg$ contains a magnetically arrested flow with field strength $B \sim 10^{2}$--$10^{4}$~G near the horizon \citep{Tchekhovskoy2011,Narayan2022} and a jet along the spin axis. Equatorial swarm orbits outside the disk body ($r \gtrsim 10^{2}\,\rg$, inclined to avoid both disk and jet) survive; structures inside $\sim 30\,\rg$ face erosion by the flow itself. We adopt $r_{\mathrm{in}} \simeq 10^{2}\,\rg$ as the practical inner boundary of a \emph{fueled} configuration, and $r_{\mathrm{ISCO}}$ for a dormant one.

\textbf{Formation coherence under relativistic precession.} A constraint with no single-star analogue: swarm elements at different radii and inclinations precess differentially. Pericenter advance and Lense--Thirring nodal precession scale as $(\rg/r)$ and $\astar(\rg/r)^{3/2}$ per orbit; at $r \sim 10^{2}\,\rg$ with orbital periods of hours, a formation spanning even a few per cent in radius decoheres in days to weeks. Rigid geometric formations (the flat rings and discs whose stability \citet{McInnes2026} and \citet{Raval2024} analyze around stars) are therefore excluded deep in the envelope; what survives is either axisymmetric shells, which are precession-invariant by symmetry, or swarms that treat element phases as free and route function through communication rather than geometry. This is the Kerr-specific amendment to the single-star stability literature: the deep envelope selects not only \emph{swarm} over \emph{monolith} (tides, above) but \emph{symmetric or phase-agnostic} swarms over shaped formations.

\subsection{Outer boundary}
\label{sec:outer}

\textbf{Cluster truncation.} The hole dominates the cluster potential inside its influence radius
\begin{equation}
r_{\mathrm{infl}} = \frac{GM}{\sigma^{2}}
\simeq 0.2\,\mathrm{pc}
\left(\frac{M}{2\times10^{4}\,\msun}\right)
\left(\frac{\sigma}{21\,\mathrm{km\,s^{-1}}}\right)^{-2}
\simeq 4\times10^{4}\,\mathrm{AU}.
\end{equation}
Orbits bound to the hole are progressively stripped by the fluctuating cluster field as $r \to r_{\mathrm{infl}}$; by analogy with tidal truncation of binaries in clusters \citep{BinneyTremaine2008,HeggieHut2003}, orbits are secure only for $r \lesssim 0.1\,r_{\mathrm{infl}} \sim 4\times10^{3}$~AU.

\textbf{Stellar flybys.} The operative outer constraint is sharper than truncation: stellar traffic. The rate at which cluster stars pass within impact parameter $b$ of the hole, including gravitational focusing, is
\begin{equation}
\Gamma(b) = n_{\star}\,\pi b^{2}\,v_{\mathrm{rel}}
\left(1 + \frac{2GM}{b\,v_{\mathrm{rel}}^{2}}\right),
\label{eq:flyby}
\end{equation}
with $n_{\star} \simeq 10^{4}\,\mathrm{pc^{-3}}$ (for mean stellar mass $0.3\,\msun$ at the quoted mass density) and $v_{\mathrm{rel}} \simeq \sqrt{2}\,\sigma \simeq 30\,\mathrm{km\,s^{-1}}$. At $b = 10^{3}$~AU the focusing term is $\sim 60$ and Eq.~\ref{eq:flyby} gives one passage per $\sim 10^{3}$~yr; at $b = 10^{2}$~AU, one per $\sim 10^{4}$~yr; at $b = 10$~AU, one per $\sim 10^{5}$~yr (focusing-dominated regime, $\Gamma \propto b$). Each passage tidally perturbs swarm orbits at the impulsive level $\Delta v / v_{\mathrm{orb}} \sim 2(m_{\star}/M)\,(a/b)^{2}\,(v_{\mathrm{orb}}/v_{\mathrm{rel}})$ for structure semi-major axis $a < b$, saturating at $\sim 2(m_{\star}/M)(v_{\mathrm{orb}}/v_{\mathrm{rel}})$ for penetrating passages, where $v_{\mathrm{orb}} = (GM/a)^{1/2}$. The mass ratio controls everything: at $m_{\star}/M \sim 10^{-5}$ even a penetrating passage delivers $\Delta v/v_{\mathrm{orb}} \lesssim 10^{-2}$, so no single flyby disrupts a hole-bound orbit, and the hazard is cumulative, a random walk in eccentricity. Two further protections apply. Direct star--element scattering requires approach within $\sim 2Gm_{\star}/(v_{\mathrm{rel}}v_{\mathrm{orb}})$, a fraction $\lesssim 10^{-7}$ of the orbital cross-section per penetrating passage: negligible. And deep in the envelope the encounter duration $b/v_{\mathrm{rel}}$ exceeds the orbital period by orders of magnitude, so the impulsive estimate above is an overestimate: adiabatic invariance suppresses the coupling of slow perturbations to fast orbits exponentially \citep{BinneyTremaine2008,HeggieHut2003}.

\subsection{Monte Carlo of flyby histories}
\label{sec:mc}

We quantify the cumulative hazard with a Monte Carlo over encounter statistics (Figure~\ref{fig:envelope}; method and parameters in Appendix~A, code in the paper repository): impact parameters from the focused distribution of Eq.~\ref{eq:flyby} within $b < 30a$, velocities from a Maxwellian at dispersion $\sigma$, masses from a two-component mass function (main sequence plus white-dwarf tail; \citealt{Haberle2024oMEGACatII}), impulsive kicks with the saturated form above, and survival defined as eccentricity random-walking to orbit-crossing ($e \sim 0.5$). The calculation is deliberately conservative in omitting adiabatic suppression.

The result is simpler than the constraint inventory led us to expect, and we flag it as a revision of our own initial estimates: \emph{flybys never set the outer boundary}. The median impulsive-diffusion lifetime is $\sim 3\times10^{9}$~yr, roughly flat from $10^{-3}$ to $4\times10^{3}$~AU, because the kick variance per encounter and the encounter rate scale inversely ($\langle\delta^{2}\rangle \propto a^{-1}$ from the saturated penetrating passages that dominate it, $\Gamma \propto a$ in the focused regime), so their product is scale-free. A conservative floor comparable to the cluster age everywhere, exponentially lengthened by adiabaticity at depth, means station-keeping against stellar traffic is a trim budget rather than a survival requirement. The operative outer boundary is therefore the cluster stripping radius of \S\ref{sec:outer}, $\sim 0.1\,r_{\mathrm{infl}} \sim 4\times10^{3}$~AU, with the flyby Monte Carlo demoted to confirming that everything interior to it is dynamically quiet.

\subsection{The envelope}
\label{sec:envelopesum}

Combining \S\ref{sec:inner}--\S\ref{sec:mc}: for the fiducial system, engineered swarms (symmetric or phase-agnostic, per the precession constraint) persist passively from $r \simeq 10^{2}\,\rg$ ($2\times10^{-2}$~AU, fueled) or $r_{\mathrm{ISCO}}$ (dormant) out to the cluster stripping radius $\sim 4\times10^{3}$~AU, with a conservative flyby-diffusion floor of $\sim 3\times10^{9}$~yr throughout and exponentially longer protected lifetimes at depth. Beyond the stripping radius, orbits bound to the hole are not durable and only cluster-orbiting architectures (outside this paper's scope) remain.

The envelope is the paper's first result (Figure~\ref{fig:envelope}). Two corollaries matter downstream. First, the thermodynamically favored location (deep, near the flow) is also the dynamically safest, doubly so once adiabatic protection is counted: the environment does not penalize the architecture Paper A's physics prefers, which is a nontrivial consistency check the hypothesis could have failed. Second, the envelope concentrates any engineered mass within $\sim 4\times10^{3}$~AU $\sim 10^{-1}\,r_{\mathrm{infl}}$ of the hole, with the thermodynamic gradient pushing occupation deep into the arcsecond-unresolvable interior: observable only through its energetic and dynamical residue.

\begin{figure}[tbp]
\centering
\includegraphics[width=0.85\linewidth]{figs/fig1_envelope.pdf}
\caption{Survival time of hole-bound swarm structures against cumulative stellar-flyby perturbation, from the impulsive Monte Carlo of \S\ref{sec:mc} (median and 16--84 per cent band; $2\times10^{4}$ histories per radius bin). The floor is conservative: adiabatic invariance suppresses the deep-envelope hazard below the impulsive estimate. Vertical markers show the inner boundaries and the cluster stripping radius; flybys never bind inside it.}
\label{fig:envelope}
\end{figure}

That is the subject of \S\ref{sec:residue}, after two remaining environmental questions: hazards other than flybys, and fuel.

\subsection{Hazards beyond flybys}
\label{sec:hazards}

Three further environmental processes deserve quantitative treatment. Each is a candidate defeater; none defeats, though the first comes closest.

\textbf{Tidal disruption events.} A cluster-core IMBH tidally disrupts stars scattered into its loss cone. Rates for IMBHs in evolved globular clusters are $\sim 10^{-8}$--$10^{-7}\,\mathrm{yr^{-1}}$ per cluster for main-sequence stars, with white-dwarf disruptions rarer by a factor of $\sim 30$--$100$ \citep{Fragione2018}. During a disruption the accretion luminosity approaches Eddington, $L \sim 3\times10^{42}\,\mathrm{erg\,s^{-1}}$ for the fiducial mass, sustained for months to years of fallback: the flux at $r = 1$~AU is $\sim 10^{9}$ times the solar constant, and no plausible material hardening survives it in place. Feasibility therefore requires that a TDE be survivable by \emph{response} rather than endurance: loss-cone stars are in principle identifiable and trackable long before disruption (the swarm sits at the bottom of the potential it would need to monitor), and the months-long fallback rise gives further warning, so temporary evacuation outward along the envelope, or shadowing at large inclination, is an engineering requirement we impose on any architecture rather than a reason to exclude one. The expected recurrence time of $\gtrsim 10^{7}$~yr then sets the natural amortization horizon of the installation: a $10^{8}$--$10^{9}$-yr occupation must plan for several such events. Two observational corollaries. The TDE duty cycle also reconciles \Heng{} with the present silence at no extra cost: post-TDE fallback outshines any steady configuration for $\sim 10^{2}$--$10^{3}$~yr per $\sim 10^{7}$-yr recurrence, a duty cycle of $\lesssim 10^{-4}$, so catching \ocen{} mid-flare was never likely under any hypothesis. And TDEs cut the other way, as an observational gift: any future TDE flare in \ocen{} both confirms the hole and activates the variability channel of \S\ref{sec:madreg} at high signal-to-noise, and the framework of \S\ref{sec:adjudication} treats a TDE with anomalous light-curve regulation as one of the few single-epoch events that can move \lnK{} substantially.

\textbf{Black-hole wander.} The hole is a Brownian particle in the stellar bath: $N$-body characterizations give r.m.s. displacements of order $10^{-2}$--$10^{-1}$~pc for $10^{3}$--$10^{4}\,\msun$ holes in \ocen-like cores, with wander amplitude scaling inversely with hole mass \citep{deVita2018}. For the fiducial $2\times10^{4}\,\msun$ hole the expected wander is $\lesssim 10^{-2}$~pc $\sim 2\times10^{3}$~AU: larger than the envelope. Structures bound to the hole simply ride along (the swarm orbits the hole, and hole plus swarm wander together through the cluster), so wander does not threaten the installation; what it threatens is the \emph{observer's} astrometry, since the kinematic center and the hole need not coincide at the $10^{-2}$-pc level. Paper C's astrometric programs already marginalize over center position; we note here that the wander amplitude is itself mass-dependent and therefore carries independent information about $M$ in long-baseline proper-motion data.

\textbf{Gas drag and erosion.} \ocen{} retains a measurable intracluster medium: sodium-absorption mapping and pulsar dispersion in comparable clusters give central ionized densities $n_{e} \sim 0.1$--$0.3\,\mathrm{cm^{-3}}$ \citep{Freire2001,Abbate2018,oMEGACatVII2026}. Ram-pressure drag on a swarm element of areal density $\Sigma \sim 1\,\mathrm{g\,cm^{-2}}$ at orbital velocity $\sim 10^{3}\,\mathrm{km\,s^{-1}}$ (the deep envelope) removes a fractional momentum $\sim \rho_{\mathrm{gas}} v t / \Sigma$ per time $t$: at the measured densities this is $\lesssim 10^{-6}$ per kyr, negligible. Sputtering and dust impacts are similarly small in an old cluster with no star formation and depleted debris populations. Gas matters for fuel, however, which is the next subsection.

\subsection{The fuel budget}
\label{sec:fuel}

The Blandford--Znajek architecture of Paper A needs mass supply, and the measured medium bounds it. Bondi accretion from gas of density $n_{e} \simeq 0.23\,\mathrm{cm^{-3}}$ \citep{oMEGACatVII2026} at relative velocity $\sim \sigma$ gives
\begin{equation}
\dot{M}_{\mathrm{B}} = \frac{4\pi (GM)^{2}\rho_{\mathrm{gas}}}{(\sigma^{2}+c_{s}^{2})^{3/2}}
\simeq 3\times10^{18}\,\mathrm{g\,s^{-1}}
\left(\frac{M}{2\times10^{4}\,\msun}\right)^{2}
\left(\frac{n_{e}}{0.23\,\mathrm{cm^{-3}}}\right)
\simeq 5\times10^{-8}\,\msun\,\mathrm{yr^{-1}},
\end{equation}
about $10^{-4}$ of the Eddington rate. Two caveats bound this from both sides. It is an upper bound insofar as the density at the hole's actual location may be below the cluster mean: the accretion non-detections are consistent with the hole occupying a locally evacuated region \citep{Mahida2026}, and the Bondi radius ($\sim GM/c_{s}^{2} \sim 2\times10^{5}$~AU) samples gas the surveys average over. It is a lower bound insofar as it ignores harvesting: stellar winds from the $\sim 10^{3}$ giants inside the influence radius, or deliberately imported mass, can raise supply by orders of magnitude at the cost of visibility. At MAD-plus-spin effective efficiencies of order unity (Paper A; \citealt{Tchekhovskoy2011}), the ambient rate alone gives extractable power $P \sim \dot{M}_{\mathrm{B}}c^{2} \sim 3\times10^{32}$~W $\sim 10^{6}\,L_{\odot}$: the ambient medium alone, with no harvesting of stellar winds or imported fuel, funds a computational budget two orders of magnitude above the waste-heat ceiling derived in \S\ref{sec:wasteheat}. The binding constraint on a present-day \ocen{} installation is therefore thermal (getting rid of entropy without detection), never supply. This ordering is a nontrivial output: for stellar-mass holes in the same environment the two constraints reverse, which is an independent reason the hypothesis family selects IMBHs. The same numbers close the loop on the accretion non-detection: the deep ATCA and JWST limits require the \emph{natural} radiative efficiency of any Bondi-fed flow to be $\lesssim 4\times10^{-3}$ \citep{Mahida2026,Chen2025JWST}, which is uncomfortable for \Hgas{} at the upper end of the mass range and what \Heng{} predicts for a flow whose output is being extracted as work rather than radiation. The adjudication in \S\ref{sec:worked} prices this observation for both hypotheses rather than letting either claim it informally.

% =====================================================================
\section{Observable residue}
\label{sec:residue}

\subsection{Waste-heat floor}
\label{sec:wasteheat}

Paper A argued that horizon entropy disposal permits computation with radiated waste heat far below the Dyson-sphere expectation. Here we make the floor quantitative. Let $P_{\mathrm{comp}}$ be the power processed by the swarm and $f_{\mathrm{sink}}$ the fraction of generated entropy delivered to the horizon rather than radiated. The radiated luminosity is $L_{\mathrm{waste}} = (1 - f_{\mathrm{sink}})\,P_{\mathrm{comp}}$, emitted from the swarm photosphere at temperature set by the emitting area. For a swarm of total area $A$ at radius $r \sim 1$~AU around a dark hole, thermal emission at $L_{\mathrm{waste}}$ peaks in the mid-infrared for $L_{\mathrm{waste}} \sim 10^{-2}$--$10^{2}\,L_{\odot}$.

The engineering limit on $f_{\mathrm{sink}}$ is transport: entropy must be carried inward, against the swarm's own power draw, by mass flux or directed radiation into the horizon. A conservative accounting (Appendix~A.3) bounds $1 - f_{\mathrm{sink}} \gtrsim 10^{-4}$ for swarm architectures within the envelope, set by the fraction of processed power unavoidably intercepted and re-thermalized by swarm elements outside the beaming solid angle. This floor is an engineering estimate rather than a first-principles derivation, and the sign of its influence is worth stating: a lower true floor (better beaming than the A.3 accounting allows) weakens the mid-infrared charge against \Heng{} and moves \lnK{} toward zero, while a higher floor tightens the $P_{\mathrm{comp}}$ ceiling and strengthens it (\S\ref{sec:discussion}). The floor is therefore
\begin{equation}
L_{\mathrm{waste}} \gtrsim 10^{-4}\,P_{\mathrm{comp}},
\end{equation}
and any mid-infrared point-source limit $L_{\mathrm{lim}}$ translates directly into $P_{\mathrm{comp}} \lesssim L_{\mathrm{lim}}/(1-f_{\mathrm{sink}}) \lesssim 10^{4}\,L_{\mathrm{lim}}$: for the sub-$L_{\odot}$ limits at the kinematic center from the JWST photometry of Paper C program 1 \citep{Chen2025JWST}, the current bound is $P_{\mathrm{comp}} \lesssim 10^{3}$--$10^{4}\,L_{\odot}$ for fiducial swarm geometry, which is far below the Eddington-scale budgets of Paper A but far above zero. An engineered system operating at the ceiling of this bound is consistent with all current data. Deeper MIR photometry tightens the bound linearly. Figure~\ref{fig:plane} assembles the waste-heat wedge, the transport floor, and the fuel ceiling of \S\ref{sec:fuel} into the constraint plane that defines the surviving \Heng{} parameter space.

\begin{figure}[tbp]
\centering
\includegraphics[width=0.85\linewidth]{figs/fig2_constraint_plane.pdf}
\caption{The $(P_{\mathrm{comp}},\,1-f_{\mathrm{sink}})$ constraint plane for an \ocen{} installation. Red: excluded by the JWST mid-infrared point-source limit. Grey hatching below the dashed grey line: the adopted entropy-transport floor $1-f_{\mathrm{sink}} = 10^{-4}$, a parameter carried over from the Appendix~A.3 accounting rather than a data-driven boundary (\S\ref{sec:discussion}). Blue dashed: the ambient Bondi fuel ceiling ($\sim 8\times10^{5}\,L_{\odot}$; imported mass moves it right at the cost of visibility). The surviving region is dormant or low-power and deep.}
\label{fig:plane}
\end{figure}

\subsection{The MAD-regulation signature}
\label{sec:madreg}

The second channel is new to this paper. General-relativistic magnetohydrodynamic simulations establish that natural accretion in the magnetically arrested state is intrinsically episodic: magnetic flux accumulates at the horizon, chokes the inflow, and erupts in quasi-regular flux-expulsion events, producing characteristic variability in jet power and radiative output with a MAD $>$ intermediate $>$ SANE hierarchy in both luminosity and its variance \citep{Tchekhovskoy2011,Narayan2022,Universe2026MAD}. The eruption cycle time is $\sim 10^{2}$--$10^{3}\,\rg/c$, which for $M = 2\times10^{4}\,\msun$ is $\sim 10$--$10^{2}$~s: minutes, comfortably within the cadence of radio and X-ray monitoring.

A system engineered for steady power extraction has an incentive absent in nature: eruptions are interruptions. Regulating the delivered mass and flux (the control variable identified by the simulations is horizon magnetic flux itself) suppresses the eruption cycle and its variability signature. We therefore forward-model the residue as a \emph{variability deficit}: an accreting IMBH whose radio/X-ray power spectrum lacks the flux-eruption band that GRMHD baselines predict for its luminosity and inferred state. Quantitatively, we define the regulation statistic $R$ as the ratio of integrated power in the eruption band ($10^{-3}$--$10^{-1}$~Hz for the fiducial mass) to the GRMHD-calibrated expectation at matched mean luminosity; natural MADs populate $R \sim 1$ with factor-of-few scatter, and $R \ll 10^{-1}$ sustained over many cycle times has no identified natural mechanism (Figure~\ref{fig:rstat})
\todo[inline]{Verify against the 2026 GRMHD variability tables before v1.0; INT-state contamination is the main confusion risk.}.

\begin{figure}[tbp]
\centering
\includegraphics[width=0.8\linewidth]{figs/fig4_r_statistic.pdf}
\caption{Schematic definition of the regulation statistic $R$: power spectra of an accreting IMBH at matched mean luminosity under natural MAD accretion (flux-eruption band populated, $R \sim 1$) and regulated inflow ($R \ll 10^{-1}$). Band location scales inversely with black-hole mass; values shown for $M = 2\times10^{4}\,\msun$.}
\label{fig:rstat}
\end{figure}

The signature has three properties valuable for adjudication. It is \emph{conditional}: it activates only if accretion is ever detected, so it presently constrains nothing (\ocen's hole is electromagnetically silent; \citealt{Mahida2026}). It is \emph{differential}: it compares a source against a physics baseline at matched parameters, canceling many systematics. And it is \emph{disprovable in place}: detection of normal flux-eruption variability from a future \ocen{} accretion flare would count \emph{against} \Heng{} in the framework below, making the channel one of the few that can move evidence in both directions.

\subsection{Relic residue: the abandoned swarm}
\label{sec:relic}

Persistence is conditional on maintenance, and \citet{Lacki2025Cascade} has shown what maintenance failure means for megastructures: once guidance fails, swarm elements collide, and a collisional cascade grinds the population to dust on a timescale of roughly the orbital period divided by the covering fraction. The deep envelope makes this dramatic. At $r \sim 1$~AU around the fiducial hole the orbital period is $\sim 3$~days; for covering fractions $10^{-4}$--$10^{-2}$ the first-collision timescale $\sim P/f_{\mathrm{cov}}$ is years to decades, and the full cascade completes within $10^{2}$--$10^{3}$~yr, orders of magnitude faster than for the stellar-orbit megaswarms Lacki considered. The end state differs too, and the difference is observationally decisive: cascade debris around a star settles into a long-lived warm dust population, a passive relic technosignature, whereas debris around a black hole is progressively drained, its periapsis distribution fed toward the loss cone by continuing collisions and flyby perturbations and its finest grindings coupled to whatever gas flow exists\todo[inline]{Drainage timescale needs a real calculation for v1.0; the qualitative asymmetry against the stellar case is robust, the rate is not yet.}. The hole removes its own debris.
% ---------------------------------------------------------------------
% SCAFFOLD for the v1.1 drainage-rate calculation (Limitations item 2).
% Physics to be checked by TS before any number replaces the 1e2-1e3 yr estimate.
% Setup, not a result:
%   1. Debris population: post-cascade size distribution dN/ds ~ s^{-3.5}
%      (collisional steady state), total mass = swarm mass, semi-major axis
%      distribution inherited from the pre-abandonment swarm (r ~ 1 AU fiducial).
%   2. Loss cone: J_lc^2 = 2 G M r_dis, with r_dis the grain-destruction radius
%      (sublimation, not tidal, for grains). Empty- vs full-loss-cone regime set
%      by comparing the per-orbit angular-momentum diffusion to J_lc.
%   3. Diffusion drivers, to be summed: (a) mutual grain collisions during the
%      ongoing cascade, (b) stellar flyby torques from the fig1 encounter rate,
%      (c) Poynting-Robertson / gas drag against the Bondi flow at n_e = 0.23 cm^-3
%      -- (c) plausibly dominates for the finest grindings and has no analogue in
%      the stellar case, which is why the asymmetry is robust independent of rate.
%   4. Deliverable: t_drain(s) as a function of grain size, then a mass-weighted
%      relic-lifetime distribution to compare against the warm-dust persistence
%      time around a star.
% --------------------------------------------------------------------- An abandoned engineered-IMBH system therefore passes through a brief bright phase (cascade grinding plus enhanced accretion luminosity as debris drains) and then reverts to a naked quiescent hole, indistinguishable from \Hgas{} except possibly through spin.

Three consequences follow. First, a fourth hypothesis joins the menu at zero structural cost: \ensuremath{H_{\mathrm{eng}}^{\mathrm{relic}}}, a formerly engineered now-abandoned system, whose observables today equal \Hgas{} plus high spin, and whose prior couples to the goal-stability open problem flagged in Paper B. Second, the searchable relic window is short, so population-level searches for dead installations (the analogue of the dust-relic searches now proposed for stellar systems) have low yield around IMBHs specifically; non-detection of relic dust in \ocen{} carries almost no evidence either way, and the framework scores it accordingly. Third, spin becomes the only durable fossil, and an objection must be met before relying on it: Blandford--Znajek extraction spins the hole \emph{down}, and MAD jets remove angular momentum faster than accretion supplies it \citep{Narayan2022}, so one might expect long use to erase the very fossil we propose reading. The rates rescue the argument: spin-down requires processing a mass of order the hole's own, and at the ambient fueling of \S\ref{sec:fuel} the timescale is $M/\dot{M}_{\mathrm{B}} \sim 4\times10^{11}$~yr. Only a civilization importing mass at far above ambient rates for $\gg 10^{9}$~yr would measurably despin the hole, and that regime is separately excluded by the waste-heat ceiling. High spin at low luminosity therefore remains the discriminating fossil: the engineered-history hypotheses predict spins drawn from the high tail (selection at arrival, negligible erosion in use), while the natural in-situ channels predict the moderate spins of chaotic hierarchical growth \citep{GonzalezPrieto2025}. The information forecast of \S\ref{sec:worked} inherits this weighting.

\subsection{Predicted-signature table}
\label{sec:nulltable}

Table~\ref{tab:residue} summarizes what each channel looks like under the engineered hypothesis and the two astrophysical nulls defined in \S\ref{sec:hypotheses}.

\begin{table}[htbp]
\centering
\small
\caption{Forward-modeled observables under the three hypotheses of \S\ref{sec:hypotheses}. ``Silent'' means below current limits. GRMHD baseline from \citet{Universe2026MAD}.}
\label{tab:residue}
\begin{tabular}{p{3.4cm}p{3.5cm}p{3.5cm}p{3.5cm}}
\toprule
Channel & \Heng{} (engineered) & \Hgas{} (quiescent IMBH) & \Hsub{} (remnant subcluster) \\
\midrule
MIR point source & floor at $10^{-4}P_{\mathrm{comp}}$; may sit near current limits & silent; $L \propto \dot{M}$ of ambient gas only & silent; no compact source \\
Radio continuum & silent while dormant; regulated (low-$R$) if fueled & silent or canonical fundamental-plane track & silent \\
Variability ($R$ statistic) & $R \ll 1$ if accreting & $R \sim 1$ if accreting & n/a \\
LISA EMRI spin & high-tail spin (selection at arrival; \S\ref{sec:relic}) & moderate spin (in-situ growth prior; \citealt{GonzalezPrieto2025}) & no EMRI (no massive central object) \\
Fast-star kinematics & IMBH-like point mass & IMBH point mass & extended mass profile \\
MSP timing & point-mass potential & point-mass potential & extended potential (current data lean this way; \citealt{BanaresHernandez2025}) \\
\bottomrule
\end{tabular}
\end{table}

The table exposes the central inferential difficulty, stated plainly in Paper A and now quantified: \emph{silence is predicted by everything}. \Heng{} and \Hgas{} are near-degenerate in every electromagnetic channel while the system is dormant. The channels that separate them are spin (engineered selection favors the high tail; in-situ growth favors moderate values), the $R$ statistic (active only during accretion episodes), and the waste-heat floor (which separates them only near the sensitivity ceiling). This degeneracy structure, rather than any single measurement, is what the adjudication framework must manage.

% =====================================================================
\section{The adjudication framework}
\label{sec:adjudication}

\subsection{Hypothesis space}
\label{sec:hypotheses}

We adjudicate among three hypotheses about the \ocen{} center:
\begin{itemize}
\item \Hgas: a quiescent IMBH of mass $10^{4}$--$5\times10^{4}\,\msun$, electromagnetically silent for want of fuel, spin drawn from an uninformative prior. The parsimonious reading of the fast stars \citep{Haberle2024Nature}.
\item \Hsub: no IMBH; the central mass is an extended subcluster of stellar remnants of $\simeq 2$--$3\times10^{5}\,\msun$ within the core, as favored by the pulsar-timing analysis \citep{BanaresHernandez2025}.
\item \Heng: an IMBH hosting engineered infrastructure within the envelope of \S\ref{sec:envelope}, with observables forward-modeled in \S\ref{sec:residue} and parameters $(P_{\mathrm{comp}}, f_{\mathrm{sink}}, R, \astar)$ assigned the priors stated in Appendix~B.
\end{itemize}
The menu is explicit and extensible; adding hypotheses (an exotic compact object, instrumental systematics on the fast stars, the abandoned-system variant \ensuremath{H_{\mathrm{eng}}^{\mathrm{relic}}} of \S\ref{sec:relic}) changes bookkeeping, never structure. We follow the convention that \Heng{} is reported \emph{only} as odds against the best-performing null, never against a strawman.

\subsection{Per-messenger Bayes factors}
\label{sec:perchannel}

For each observational channel $i$ with data $d_{i}$, the Bayes factor between hypotheses $H_{a}, H_{b}$ is
\begin{equation}
K_{i}^{ab} = \frac{\int \mathcal{L}(d_{i}\,|\,\theta_{a})\,\pi(\theta_{a})\,d\theta_{a}}
{\int \mathcal{L}(d_{i}\,|\,\theta_{b})\,\pi(\theta_{b})\,d\theta_{b}},
\end{equation}
with channel likelihoods built from the forward models of \S\ref{sec:residue} and, following radio-SETI practice, an explicit interference/systematics component mixed into every likelihood \citep{Sheikh2021,Luan2026} so that ``instrumental artifact'' is priced inside each channel rather than adjudicated informally afterward. Priors on \Heng{} parameters are pre-registered (Appendix~B) and varied over stated ranges in sensitivity analysis; the framework's outputs are always reported as (odds, prior-sensitivity band) pairs, never as bare odds.

\subsection{Cross-messenger combination}
\label{sec:combination}

Channels are combined multiplicatively where independent, $K_{\mathrm{tot}} = \prod_{i} K_{i}$, with two corrections imported from gravitational-wave counterpart methodology \citep{Ashton2018,Breschi2024}. First, a \emph{coincidence term}: hypotheses that predict correlated anomalies across channels (as \Heng{} does for MIR excess and low $R$ during an accretion episode) earn a likelihood contribution from the observed coincidence structure itself, computed from the joint forward model rather than the product of marginals. Second, an \emph{independence audit}: channels sharing calibrators, atmospheric paths, or reduction pipelines are grouped and their shared systematics marginalized jointly before multiplication. Both corrections are conservative in effect: the coincidence term rewards only pre-registered correlation patterns, and the audit only ever weakens evidence.

\subsection{Decision thresholds}
\label{sec:thresholds}

We adopt the logarithmic odds scale of \citet{Kass1995} with pre-registered action bands, stated here for \lnK{} of \Heng{} over the best null:
\begin{center}
\small
\begin{tabular}{lll}
\toprule
Band & \lnK & Action \\
\midrule
Null-favored & $< 0$ & routine monitoring; \Heng{} parameter space shrinks \\
Uninformative & $0$ to $1$ & no action; report in campaign updates \\
Anomaly & $1$ to $3$ & targeted follow-up on the discriminating channels \\
Strong anomaly & $3$ to $5$ & independent-team replication; data release \\
Candidate & $> 5$ sustained across $\geq 2$ messengers & community adjudication per post-detection protocols \\
\bottomrule
\end{tabular}
\end{center}
Two design points. The ``candidate'' band requires multi-messenger support by construction, honoring Paper C's two-messenger rule; no single channel, at any significance, can reach it, because single-channel anomalies are where every historical false alarm has lived \citep{Sheikh2021}. And the kill conditions of Paper A map onto the same scale in the other direction: the pre-registered falsifiers (e.g., LISA measurement of low spin; resolution of the mass tension in favor of \Hsub) enter as ordinary likelihood terms and drive \lnK{} negative, so confirmation and falsification run through one pipeline rather than two standards. The framework is deliberately instrument-agnostic: nothing in \S\ref{sec:perchannel}--\S\ref{sec:thresholds} references \ocen, and the machinery applies unchanged to any technosignature target with a defined null menu, complementing the qualitative Rio scale \citep{Forgan2018} with a quantitative interior.

% =====================================================================
\section{Worked example: Omega Centauri today}
\label{sec:worked}

We now run the current \ocen{} data through the framework. Inputs: the fast-star kinematics \citep{Haberle2024Nature}, the $N$-body modeling \citep{GonzalezPrieto2025}, the MSP timing bound \citep{BanaresHernandez2025} with the TRAPUM 2026 mass limit, the deep radio silence \citep{Mahida2026,Tremou2018}, and the JWST infrared limits \citep{Chen2025JWST}. The computation implements the Appendix-B model (deliberately minimal: hard-threshold MIR likelihood, log-flat priors, $2\times10^{6}$ prior draws); code in the paper repository, results in Figure~\ref{fig:lnk}.

\begin{itemize}
\item \textbf{\Hgas{} vs \Hsub:} the interesting contest, and the framework's finding is that it is currently a stand-off dominated by the kinematics-versus-timing tension: the fast stars pull toward a point mass, the pulsars toward an extended one, and the combined $|\lnK|$ is $\lesssim 1$ with sign depending on how the two datasets' systematics are weighted. As an independent check, the \texttt{imbh-constraints} aggregation engine of the series' public toolchain, run on its nine-constraint curated compilation, returns an empty jointly-allowed mass window (formal tension verdict: the kinematic floor exceeds the tightest model-dependent ceiling), confirming that no point-mass value satisfies all published constraints at face value. The framework adds discipline, and no verdict, to a tension the community already knows; per series policy the tension is recorded, never collapsed.
\item \textbf{\Heng{} vs best null:} \lnK{} sits in the null-favored band, as it should. The only active channel today is waste heat, and it yields $\lnK = -0.40$ at fiducial priors (prior-sensitivity band $-0.90$ to $-0.07$ as the dormancy prior runs from $0.1$ to $0.9$): the electromagnetic silence sometimes informally cited as consistent with engineering is, in the accounting, mildly \emph{against} \Heng, because the JWST limit excludes the $\sim 66$ per cent of the active-installation prior volume with $L_{\mathrm{waste}}$ above it, while every other channel is inactive or degenerate and contributes nothing. The leak prior's lower bound is an engineering estimate (\S\ref{sec:wasteheat}, Appendix~A.3): repeating the computation with the prior extended down to $1-f_{\mathrm{sink}} = 10^{-6}$ gives $\lnK = -0.28$ instead of $-0.40$, so a softer transport bound softens the conclusion without reversing it. The surviving configurations are dormant or low-power, deep-envelope systems (Figure~\ref{fig:plane}). Quantifying that shrinkage is the point; the hypothesis pays for its unfalsified survival in measured parameter volume.
\item \textbf{Information forecast:} ranking the Paper C programs by expected $|\Delta\lnK|$ per unit cost, the leaders are (1) resolution of the mass tension by continued MSP timing plus Gaia DR4 astrometry, which arbitrates \Hgas{} vs \Hsub{} and thereby moves \Heng's denominator; (2) any detection of accretion at any level, which activates the $R$ statistic, the only cheap channel with large discriminating power in both directions; (3) the LISA-era spin measurement, the single largest mover but the least schedulable: intermediate-mass-ratio inspiral rates per IMBH are $\sim 10^{-9}$--$10^{-6}\,\mathrm{yr^{-1}}$ \citep{Mandel2008,AmaroSeoane2018,Fragione2018}, so \ocen{} yielding its own inspiral during the mission is unlikely and the spin constraint will more probably arrive statistically, from the population of \ocen-like systems LISA does catch, than from \ocen{} itself. The relic analysis of \S\ref{sec:relic} raises the stakes on this channel: spin is the one fossil that survives abandonment, so the LISA-era population measurement adjudicates the engineered-history hypotheses for the whole Galactic IMBH population at once, and the per-target frameworks of this paper compose naturally into that population-level test. Deeper radio silence, by contrast, now moves \lnK{} weakly for every pairing: for this hypothesis menu the existing limits are already deep enough that further depth buys little discrimination.
\end{itemize}

\begin{figure}[tbp]
\centering
\includegraphics[width=0.85\linewidth]{figs/fig3_lnk.pdf}
\caption{Per-channel $\lnK$ (\Heng{} vs \Hgas) on 2026 data, from the Appendix-B computation. Waste heat is the only live channel; its $-0.40$ (whiskers: dormancy-prior sensitivity; extending the leak prior down to $10^{-6}$ moves the central value to $-0.28$) is the entire present evidence budget. The empty bars are the framework's forecastable headroom: each activates with the observations ranked in the information forecast.}
\label{fig:lnk}
\end{figure}

% =====================================================================
\section{Discussion}
\label{sec:discussion}

\subsection{Limitations}

The feasibility envelope depends on cluster-center parameters that carry real uncertainty (mass function of the remnant tail, central density profile inside the influence radius) and on a material-strength scale taken from present technology; both are varied in Appendix~A, and the envelope boundaries move by factors of a few, not orders of magnitude. The MAD-regulation channel inherits the systematic uncertainties of the GRMHD baseline literature, which is simulation-calibrated rather than observationally calibrated at IMBH masses. The adjudication framework shares the standard vulnerability of Bayesian model selection to unmodeled hypotheses: it ranks the menu it is given, and a surprise outside the menu (an astrophysical mechanism not yet imagined) would be mis-scored until added. The mitigation is procedural rather than mathematical: the menu is public, extensible, and versioned.

Two numbers used above deserve explicit qualification at the claim level. \textbf{The transport floor.} The bound $1 - f_{\mathrm{sink}} \gtrsim 10^{-4}$ of \S\ref{sec:wasteheat} is an order-of-magnitude engineering accounting rather than a derivation from first principles. A real derivation needs a radiator geometry for the disposal channel and a beaming solid-angle model for re-interception within the swarm, replacing the covering-fraction argument of Appendix~A.3 with a radiative-transfer calculation. Its influence on the adjudication is bounded and signed: extending the leak prior down to $10^{-6}$ moves \lnK{} from $-0.40$ to $-0.28$, well inside the dormancy-prior band of $[-0.90, -0.07]$, so no conclusion of \S\ref{sec:worked} turns on the exact value. \textbf{The drainage timescale.} The $10^{2}$--$10^{3}$~yr figure of \S\ref{sec:relic} is an estimate. The asymmetry it rests on is robust, since relic dust persists around a star while debris around a hole is fed toward the loss cone and removed, but the rate is not. Closing that gap requires a loss-cone refill calculation for the debris population, tracking collisional periapsis diffusion and the coupling of the finest grindings to the ambient flow; we state the gap rather than supply a number that calculation has not yet produced.

\subsection{Relation to the series and beyond}

Within the series, this paper closes the constructive quadrant: A argued the destination is attractive, and \S\ref{sec:envelope} finds the destination is habitable by infrastructure, with the thermodynamically preferred region also the dynamically safest. It also arms Paper C's campaign with the scoring machinery its two-messenger policy presupposed, and returns to Paper D a refined survival input ($p_{s}$ now decomposable into transit and residence terms, the latter bounded here). Beyond the series, the two exportable products are the envelope method, applicable to any proposed megastructure environment with a stated perturber population, and the adjudication framework, applicable to any technosignature program willing to pre-register its nulls.

\subsection{Conclusion}

The engineered-IMBH hypothesis survives its first engineering audit and its first quantitative adjudication, in both cases by narrowing: infrastructure persists only within $\sim 4\times10^{3}$~AU of the hole, and the surviving parameter space after current data is dormant, low-power, and deep. Every forthcoming measurement listed in \S\ref{sec:worked} shrinks it further or breaks it open. Either outcome is progress that a hypothesis without an envelope, a residue model, and a scoring rule could not deliver.

% =====================================================================
\appendix

\section{Flyby Monte Carlo: method and parameters}
\label{app:mc}

\textbf{Method} (implemented in \texttt{figs/fig1\_envelope.py}; fixed seed 20260717). For each of 40 logarithmic radius bins $a \in [10^{-3}, 4\times10^{3}]$~AU: (1) $2\times10^{5}$ encounters are drawn with impact parameter from the gravitationally focused cumulative distribution $Q(b) \propto v_{\mathrm{rel}}^{2}b^{2} + 2GMb$ truncated at $b_{\max} = 30a$ (contributions beyond carry $<10^{-3}$ of the kick variance), speed from a Maxwellian at $\sigma = 21\,\mathrm{km\,s^{-1}}$, and mass from a two-component function ($0.35\,\msun$ with weight $0.7$; $0.6\,\msun$ white-dwarf tail with weight $0.3$). (2) Each encounter contributes an impulsive eccentricity kick $\delta = 2(m_{\star}/M)\min[1,(a/b)^{2}](v_{\mathrm{orb}}/v)$; no adiabatic suppression is applied (conservative). (3) Per history ($2\times10^{4}$ per bin), survival time is the random-walk first-passage to $e = 0.5$: encounters to threshold $n_{\star} = e^{2}/\langle\delta^{2}\rangle$ with CLT scatter, divided by the total encounter rate $\Gamma(<b_{\max})$, capped at 12~Gyr. Single-passage disruption is impossible at these mass ratios ($\delta \lesssim 10^{-2}$ even for penetrating passages); direct star--element scattering requires approach within $\sim 2Gm_{\star}/(v_{\mathrm{rel}}v_{\mathrm{orb}})$ and contributes negligibly at any realistic covering fraction.

\textbf{Results.} Median lifetimes $3.4\times10^{9}$, $3.3\times10^{9}$, and $2.5\times10^{9}$~yr at $a = 10$, $10^{2}$, $10^{3}$~AU: flat, by the scale cancellation $\langle\delta^{2}\rangle \propto a^{-1}$, $\Gamma \propto a$ discussed in \S\ref{sec:mc}. Varying the mean stellar mass by $\times 2$, the density by $\times 3$, and $e_{\mathrm{cross}}$ over $0.3$--$0.7$ moves the floor by factors of a few; no variation brings any part of the envelope below $10^{8}$~yr.

\textbf{A.3: transport bound.} The floor $1 - f_{\mathrm{sink}} \gtrsim 10^{-4}$ follows from intercepted-power geometry: entropy export to the horizon requires beaming waste flux into the solid angle subtended by the disposal channel, and for swarm geometries within the envelope, the fraction of beamed power intercepted and re-thermalized by other swarm elements is bounded below by the swarm covering fraction as seen from a typical element, $\gtrsim 10^{-4}$ for the covering fractions that make the swarm computationally worthwhile.
\todo[inline]{Weakest number in the paper: a geometric estimate rather than a derivation. v1.1 should treat radiative transfer within the swarm properly.}

\section{Priors, likelihoods, and computation}
\label{app:bayes}

Implemented in \texttt{figs/fig3\_lnk.py} ($2\times10^{6}$ prior draws; fixed seed). \Heng{} parameters: $P_{\mathrm{comp}}$ log-uniform on $[1, P_{\mathrm{fuel}}]\,L_{\odot}$ with $P_{\mathrm{fuel}} = 7.6\times10^{5}\,L_{\odot}$ (the ambient Bondi ceiling of \S\ref{sec:fuel}); $1 - f_{\mathrm{sink}}$ log-uniform on $[10^{-4}, 1]$; dormancy probability $f_{d} = 0.5$, varied over $[0.1, 0.9]$ for the sensitivity band. MIR channel likelihood: hard detection threshold at $L_{\mathrm{lim}} = 1\,L_{\odot}$ (a soft threshold changes $\lnK$ by $<0.1$); data = non-detection; $P(\mathrm{quiet}\,|\,\mathrm{active}) = 0.34$, giving $\lnK_{\mathrm{MIR}} = \ln[f_{d} + (1-f_{d})\times0.34] = -0.40$ at fiducial $f_{d}$, band $[-0.90, -0.07]$. Transport-floor sensitivity (\S\ref{sec:discussion}): re-running with the leak prior log-uniform on $[10^{-6}, 1]$ raises $P(\mathrm{quiet}\,|\,\mathrm{active})$ to $0.51$ and gives $\lnK_{\mathrm{MIR}} = -0.28$, band $[-0.58, -0.05]$; the script takes the floor as the \texttt{-{}-leak-floor-dex} argument. Radio, $R$, spin, kinematic, and timing channels are inactive or degenerate between \Heng{} and \Hgas{} on current data and contribute $\lnK = 0$ each, as Table~\ref{tab:residue} requires. The \Hgas--\Hsub{} contest is checked against the \texttt{imbh-constraints} v1.0.0 aggregation engine (nine curated constraints; empty joint window, tension verdict), whose mass-window math is anchored by the series' replayed golden and CI parity gate.

\section{Reproducibility}
\label{app:repro}

Figure scripts and outputs live beside this manuscript (\texttt{paper/figs/}: \texttt{common.py} for shared constants, one script per figure, \texttt{fig1\_results.json} for machine-readable Monte Carlo output). The measurement compilation and aggregation library are the public series toolchain (\texttt{imbh-constraints} v1.0.0, Zenodo DOI 10.5281/zenodo.20689279, ASCL submitted); interactive calculator versions of the Figure~\ref{fig:envelope} and \S\ref{sec:fuel} computations are planned for omegacentauri.me. All random draws use fixed seeds; all quoted numbers regenerate from the scripts.

% =====================================================================
\section*{Acknowledgements and disclosure}

The author thanks the maintainers of the NASA Astrophysics Data System and arXiv, on which the citation verification for this work relied. \textbf{AI assistance disclosure:} drafting, citation verification, derivation checking, and figure preparation for this manuscript were performed with substantial assistance from a large language model (Claude, Anthropic), under the author's direction; the author reviewed and takes full responsibility for all claims, derivations, and references. Interactive calculators implementing the quantitative material in Sections~\ref{sec:envelope}--\ref{sec:adjudication} are available as supplementary material at \url{https://omegacentauri.me}.

\section*{Data availability}

No new observational data were generated. All cited measurements are available in the referenced publications. The figure scripts, shared constants, and machine-readable Monte Carlo output described in Appendix~\ref{app:repro} are mirrored publicly at \url{https://omegacentauri.me/papers/source/} and \url{https://omegacentauri.me/papers/figs/}, alongside the interactive calculator versions of the same computations.

% =====================================================================
\bibliographystyle{apalike}
\bibliography{references,engineered-extra}

\end{document}
